\section{Course Representations}
\label{sec:representations}

A course representation is not a neutral implementation detail. It determines which
constraints can be generated, checked, exchanged, and evaluated without reconstruction
assumptions. We therefore separate three levels. The \emph{course object} is the
task-defining spatial and traversal specification introduced in
\cref{sec:scope-definitions}: a corridor and its admissible route, a directed gate
chain, a waypoint relation, or a buoy course. A \emph{representation} encodes that
object. A \emph{concrete realization} is a particular map, mesh, world file, raster,
or simulator asset built from an encoding. This separation is practical rather than
ontological. The same closed circuit may be represented by a curvature-continuous
centerline and widths, instantiated as paired boundaries and cones, and rendered as
a textured collision mesh. None of those transformations is automatically lossless.

The evidentiary rule follows from this distinction: call a form a course
representation only when the source itself defines, emits, consumes, serializes,
inspects, or releases it as course state. A rendered mesh, physics collider,
occupancy image, coordinate buffer, or imported world is otherwise evidence about a
realization, not about the source's course model. This rule avoids a common category
error in which a simulator's internal geometry is attributed to a generator that
actually specifies only a route or an ordered gate set. It also keeps standards and
fixed benchmarks in their proper role: CommonRoad, for example, packages scenarios
with goals and constraints in a portable XML form, but the existence of this
representation does not establish a course-synthesis method
\cite{Althoff2017CommonRoadComposable}.

For future taxonomies, reports should record two fields rather than force this
distinction into a single label: (i) the \emph{primary course representation}, namely
the source-native state on which course-defining operations act, and (ii) \emph{derived
realization tags}, such as sampled polyline, raster, triangle mesh, SDF world, or
simulator scene. A tag is promoted to the primary field only when it carries the
definition needed to reconstruct the relevant course object. The families below
organize analytic differences; they do not assign a rank.

\subsection{Structural and parametric representations}

\paragraph{Topology and segment grammars.}
A grammar represents a course as a composition of typed elements and their connection
rules. Its natural state includes an ordered or directed sequence, attachment ports,
and parameters for each element; closure and branching are graph properties, not
visual consequences. PGDrive illustrates this approach by sampling elementary road
blocks and connecting their sockets to form road networks \cite{Li2020ImprovingGeneralization}.
Such a representation retains topology and local admissibility but not necessarily continuous geometry or a surface. A node--edge network supports route reasoning, but route, direction, width, and termination remain extra course state. SUMO netgenerate constructs abstract networks from grid, spider, or random controls, including edge and junction parameters
\cite{EclipseSUMONodateNetgenerate}.

\paragraph{Curves, centerlines, and boundaries.}
Parametric curves make geometry continuous and differentiable where required. A
Bezier control-point sequence can define a road curve directly; moving one control
point changes a global curve and makes curvature, self-intersection, and clearance
testable in the same parameter domain. Kl\"uck \emph{et al.} use precisely this
relationship between control points, Bezier roads, and candidate validity checks
\cite{Kluck2021AutomaticGeneration}. Track generators have also combined a segment
description with a continuous curve representation, showing that these are often
complementary stages rather than competing answers to one taxonomy question
\cite{Loiacono2011AutomaticTrack}.

A sampled centerline $c(s)$ with left and right widths $w_L(s),w_R(s)$ is compact when order, arc-length convention, orientation, and frame are declared. It retains progress coordinate and tangents, but generates boundaries only under stated surface assumptions. The TUM race-track database releases centerline coordinates with left and right widths, while treating optimized racing lines as a distinct derived product
\cite{TUMFTM2020RaceTrack}. A centerline alone loses asymmetric width, obstacle
placement, crossfall, and any nontrivial surface geometry. Conversely, a boundary
pair retains corridor extent and can express asymmetric or locally varying width, but
does not by itself declare traversal order, a centerline parameterization, or a route
at intersections. Lanelet2 makes the distinction explicit: each lanelet has one left
and one right boundary plus regulatory and neighboring relations; successor and
opposite-direction semantics cannot safely be recovered from a polygon alone
\cite{Poggenhans2018Lanelet2High}.

\paragraph{Samples, rasters, and meshes.}
Polylines, samples, occupancy rasters, signed-distance fields, and meshes suit different algorithms. Sampling makes collision queries approximate; spacing and closure conventions control error. A raster quantizes geometry, leaving units, orientation, connectivity, and sub-cell clearance contingent on resolution and a world transform. A mesh retains a collision-ready surface but usually omits route coordinate, boundary correspondence, and task progression; it can also introduce seams, holes, or normal conventions absent from primary state.

These forms may nevertheless be source-native. Isaac Lab documents terrains as meshes
obtained from height fields or mesh construction, with different flexibility and
memory costs \cite{IsaacLabNodateIsaacLab}. In that setting, terrain state is part of
the task definition rather than merely a rendering byproduct. The correct conclusion
is not that every terrain mesh is a course representation, but that the source's
declared terrain field, mesh, and task relation determine which one is primary. A
safe export records sampling density, interpolation, boundary policy, collision
geometry, frame, units, and the error or semantic loss induced by conversion.

\subsection{Task-element and world representations}

\paragraph{Gate poses and waypoint relations.}
For aerial and maritime tasks, a point set is often insufficient. A gate pose supplies location and orientation; the course also needs allowed crossing direction, tolerance or aperture, and precedence. Flightmare shows the simulator-side importance of static gates, while task semantics must still state a valid crossing \cite{Song2021FlightmareFlexible}. An ordered waypoint sequence is compact for one route but loses alternatives and branch conditions. A directed waypoint graph retains successors and route choice when edge direction and terminal states are explicit. OpenSCENARIO represents routes through waypoints and route strategies rather than treating a geographic polyline as complete semantics
\cite{ASAM2024ASAMOpenSCENARIO}.

Buoy and marker courses expose the same issue in a more concrete form. The frozen VRX
navigation asset contains explicitly posed port and starboard boundary buoys and
obstacles in an SDF world \cite{Bingham2019TowardMaritime}. That SDF realization
retains concrete object poses and models, but not necessarily an authoritative ordered
passage relation; inferring a gate sequence from color and proximity would silently
add task semantics. Gate poses, waypoint sequences, and buoy placements should
therefore carry an explicit order, partial order, or transition graph whenever
completion depends on it.

\paragraph{Graphs, maps, and hybrid world state.}
Road and map representations commonly combine topology with geometry. HDMapGen uses a
coarse-to-fine hierarchical graph, from which intersections and lanes can be rendered
\cite{Mi2021HDMapGenHierarchical}. Such graphs retain connectivity and can attach
lane-level attributes, but their geometry is only as precise as the node, edge, and
curve attributes they carry. At the other extreme, a simulator-native hybrid may
combine an SDF/scene graph, meshes, material and collision assets, traffic rules,
spawns, sensors, and a route. This is often the most executable realization, but it
couples the course to simulator semantics and can bury the task-defining subset among
visual and operational metadata. The representation should consequently identify the
minimal course subgraph and preserve a link to the enclosing world artifact, rather
than exporting an opaque scene as if it were a portable course specification.

\begin{table}[t]
  \centering
  \caption{Representation families are distinguished by the course information they
  preserve, not by their rendered appearance. ``Loss'' means information that must be
  supplied by an explicit convention or auxiliary field to reconstruct the course
  object.}
  \label{tab:representation-information}
  \small
  \begin{tabularx}{\linewidth}{@{}p{0.22\linewidth}Y Y@{}}
    \toprule
    Primary representation & Retains directly & Typical loss without auxiliary state \\
    \midrule
    Segment grammar / topology graph & Element types, connectivity, ports, construction order & Metric geometry, widths, surface, route semantics if edges are untyped \\
    Parametric or sampled centerline + widths & Longitudinal order, tangent, width profile, compact metric geometry & Surface, obstacles, junction choices, exact boundary correspondence under complex geometry \\
    Boundary pair / lanelet & Traversable extent, asymmetric bounds, local adjacency & Canonical progress coordinate, route order, elevation and rule semantics unless supplied \\
    Gate poses / waypoint graph & Pose or point state; declared sequence or successors when encoded & Corridor extent, clearance field, visual/collision geometry \\
    Raster, field, or mesh & Discretized occupancy, elevation, or collision/render surface & Exact topology, orientation, units/frame, progression, and semantic elements unless encoded \\
    Simulator-native hybrid & Executable world state and simulator bindings & Portability and a separable, simulator-independent course object \\
    \bottomrule
  \end{tabularx}
\end{table}

\subsection{Conversion, feasibility, and reporting}

Conversion is a partial, typed transformation, not a file-format detail. A centerline-to-mesh conversion can preserve width within tolerance while losing route alternatives; mesh-to-raster can preserve occupancy while changing clearance and connectivity; world export can retain poses while dropping checkpoint direction, rules, or identifiers. Closure is fragile: duplicate-endpoint, periodic, and wrapped-index conventions affect curvature, collision, and progress. Reversing vertex order can invert a normal, gate direction, or the meaning of left and right.

Simulator feasibility is consequently a property of the representation--converter--
simulator chain, not of an abstract curve alone. A usable import must declare coordinate
frame and handedness, units, gravity and elevation conventions where relevant,
collision-versus-visual geometry, asset dependencies, spawn and reset states, and the
mapping from course semantics to simulator termination or progress logic. The
CourseSpec and course-card contract proposed in \cref{sec:benchmark-protocol} provides
the appropriate place to record these fields, including semantic loss, conformance,
and raw-to-derived lineage. A successful render or import is necessary evidence of
interoperability, but it does not demonstrate preservation of topology, traversal
semantics, or dynamics-conditioned feasibility.

Accordingly, every release should preserve the source-native primary representation,
its content hash and units/frame metadata, and each derived realization with a declared
conversion contract. Evaluation can then compare courses in the representation in
which a descriptor is meaningful: graph distance for successor structure, arc-length
and curvature for centerlines, pose and order checks for gates, and geometry-aware
distance for surfaces. This preserves the distinction required by
\cref{sec:metrics}: geometry may be valid, a realization may import, and a vehicle
may still be unable to complete the declared task. Treating all three as a single
``track format'' obscures precisely the failures that a reproducible course generator
must expose.
